A diffusion language model builds a solution by filling masked tokens over several steps. We evaluate five models on GSM8K math problems and ask how those steps affect the benefit of choosing among solutions. We compare taking the first solution, voting for the most frequent final answer, and reranking solutions using next-token probabilities from an autoregressive (AR) model.

For Dream and LLaDA, voting closes a growing share of the gap to the best possible choice from 16 to 64 steps. On matched questions from 64 to 256 steps, that share falls from 69\% to 53\% on Dream and from 68\% to 40\% on LLaDA. At 256 steps AR reranking trails voting by 8.6 points on Dream. The 2.6-point gap on LLaDA has a confidence interval that includes zero. Scoring only through the final answer raises Dream's reranking accuracy by 13 points and changes the apparent link between similarity to the base model and reranking gains.
